\documentclass[11pt]{article}
\usepackage[margin=1in]{geometry}
\usepackage{amsmath,amssymb}
\usepackage{graphicx}
\usepackage{booktabs}
\usepackage{natbib}
\usepackage{hyperref}
\usepackage{authblk}
\usepackage{titlesec}
\usepackage{enumitem}

\titleformat{\section}{\normalfont\Large\bfseries}{\thesection}{1em}{}
\titleformat{\subsection}{\normalfont\large\bfseries}{\thesubsection}{1em}{}
\titleformat{\subsubsection}{\normalfont\bfseries}{\thesubsubsection}{1em}{}

\title{\textbf{An Open-Source End-to-End Pipeline for Multi-Planet Detection and Bayesian Characterisation in Kepler Long-Cadence Photometry: Validation on Real Kepler-11 Photometry}}
\author[1]{Sree Narasimha Badari Narayana Konda}
\affil[1]{Independent Researcher, India}
\date{}

\begin{document}
\maketitle

\begin{abstract}
\noindent
We present an open-source, end-to-end pipeline for the detection and Bayesian characterisation of multiple transiting exoplanets in Kepler long-cadence photometry, and validate it on genuine MAST-hosted Kepler-11 photometry (KIC~6541920, Quarters~1--17, 58{,}785 cleaned long-cadence cadences spanning a 1459.5-day baseline). Using the \texttt{transitleastsquares} (TLS) algorithm with literature-informed targeted period-search windows, the pipeline recovers all six known planets to sub-0.005\% period accuracy, except Kepler-11f, initially recovered at a 2.81\% period offset that we first (incorrectly) attributed to unmodelled transit-timing variations. A subsequent full blind search over the entire 0.5--300 day grid -- performed with no period given in advance, closing the paper's own previously-flagged targeted-vs-blind gap -- recovers Kepler-11f at 0.0066\% accuracy and SDE~$=192$, showing the original offset was a targeted-search artefact, not astrophysics; we report this correction explicitly rather than silently revising it away. The same blind search recovers five of six planets at very high significance (SDE~78--226); the sixth (Kepler-11b, the shallowest at 310~ppm) is not recovered blind, consistent with the look-elsewhere effect rather than a pipeline defect. Recovered transit depths for four of the six planets (b, c, d, e) are systematically shallower than the discovery-paper values, most pronounced for Kepler-11e (36--40\%); Kepler-11g shows no such shortfall (slightly deep, if anything), a distinction that itself weakens simple per-star explanations. Aperture dilution and revised stellar parameters were both tested directly against real MAST header data and ruled out as the cause, which remains open. A joint four-planet Metropolis--Hastings fit, run after correcting a step-size bug that produced a 0\% acceptance rate on the real, high-information dataset, converges to healthy 31--45\% acceptance and yields period precisions of order tens of seconds -- roughly $50\times$ tighter than the pipeline's own synthetic-data validation, reflecting the far larger number of observed transits (45--142) in the real 4-year baseline. In building this validation we identified and fixed two functional defects in the original codebase: a custom period-search statistic that returned a high-confidence, non-physical false detection on real data, and an MCMC proposal-scale heuristic that froze the sampler entirely on the full real dataset. We report these fixes, the depth-shortfall investigation, and the blind-search follow-up in full, including a claim we made and then had to retract when better evidence arrived.
\end{abstract}

\noindent\textbf{Keywords:} methods: data analysis --- methods: statistical --- planets and satellites: detection --- techniques: photometric --- stars: individual (Kepler-11)

\section{Introduction}

The discovery of thousands of exoplanets by the Kepler Space Telescope \citep{Borucki2010} has transformed our understanding of planetary system architectures. Kepler observed $\sim$150{,}000 stars nearly continuously for four years at 29.4-minute cadence, producing light curves of unprecedented photometric precision, enabling the transit method to reach Earth-radius sensitivity around bright main-sequence targets \citep{Koch2010}.

The transit method detects planets by the periodic dimming of a host star as an orbiting planet crosses the stellar disc. For a planet of radius $R_p$ orbiting a star of radius $R_\star$, the fractional depth of the transit is
\begin{equation}
\delta = \left(\frac{R_p}{R_\star}\right)^2 ,
\label{eq:depth}
\end{equation}
and the orbital period $P$ and semi-major axis $a$ are related through Kepler's third law:
\begin{equation}
a = \left(\frac{P^2 G M_\star}{4\pi^2}\right)^{1/3},
\label{eq:kepler3}
\end{equation}
where $M_\star$ is the stellar mass and $G$ is the gravitational constant. The dominant period-search approach in the literature has been Box Least Squares \citep[BLS;][]{Kovacs2002}; Transit Least Squares \citep[TLS;][]{Hippke2019} improves on it with a physically motivated trapezoidal template and a log-uniform period grid, achieving $\sim$10\% higher completeness at equivalent false-alarm rates.

Multi-planet systems -- which constitute $\sim$20\% of all Kepler host stars \citep{Fabrycky2014} -- present additional detection and characterisation challenges, particularly where mutual gravitational interactions induce transit timing variations (TTVs) that can bias single-period fits \citep{Lissauer2011}. This paper presents a modular, open-source pipeline spanning preprocessing, period search, catalog cross-matching, and joint Bayesian fitting, and -- the focus of this revision -- reports its validation against genuine Kepler photometry rather than a synthetic benchmark. Moving from synthetic to real data surfaced two functional defects in the original implementation and a set of real-data effects (instrumental discontinuities, TTVs, aperture dilution) that a synthetic-only validation could not have exposed. We report the pipeline, the defects found and fixed, and the resulting honest real-data performance.

\subsection{Kepler-11: A Benchmark Compact Multi-Planet System}

Kepler-11 (KIC~6541920) is a G-type dwarf ($T_{\rm eff} = 5680$~K, $M_\star = 0.961\,M_\odot$, $R_\star = 1.065\,R_\odot$; \citealt{Lissauer2011}) hosting at least six transiting planets in compact, near-coplanar orbits with periods $10.3$--$118.4$~days \citep{Lissauer2013}. Transit depths span roughly 300--1630~ppm; all six planets have independent confirmation; and significant TTVs from mutual gravitational interaction make the system a standard benchmark for both detection sensitivity and dynamical characterisation. We adopt the \citet{Lissauer2011,Lissauer2013} parameters as the literature comparison values throughout.

\begin{table}[h]
\centering
\caption{Literature Kepler-11 system parameters \citep{Lissauer2011,Lissauer2013}, used throughout as the comparison ground truth for real-data recovery.}
\label{tab:literature}
\begin{tabular}{lcccccc}
\toprule
Planet & Period (d) & $t_0$ (d) & Depth (ppm) & Duration (hr) & Radius ($R_\oplus$) \\
\midrule
Kepler-11b & 10.304 & 2.10 & 310 & 2.16 & 2.05 \\
Kepler-11c & 13.025 & 5.30 & 790 & 2.40 & 3.27 \\
Kepler-11d & 22.687 & 1.80 & 940 & 3.12 & 3.56 \\
Kepler-11e & 31.996 & 7.00 & 1630 & 3.60 & 4.69 \\
Kepler-11f & 46.689 & 3.50 & 540 & 2.88 & 2.70 \\
Kepler-11g & 118.38 & 12.00 & 890 & 4.80 & 3.44 \\
\bottomrule
\end{tabular}
\end{table}

\section{Data and Methods}

\subsection{Light Curve Source (Revised: Real MAST Photometry)}
\label{sec:data}

The pipeline was previously validated only on a synthetic light curve. This revision replaces that benchmark with genuine Kepler-11 long-cadence photometry, downloaded via the \texttt{lightkurve} interface to the MAST archive \citep{Lightkurve2018}: all 17 available long-cadence quarters (Q1--Q17) for KIC~6541920, totalling a 1459.5-day baseline closely matching the mission's nominal 4-year span.

\subsubsection{Detrending and Outlier Rejection (Revised: Per-Quarter Detrending Required)}

Raw Kepler light curves contain long-term trends from stellar variability, spacecraft systematics, and instrumental contamination that must be removed without attenuating the transit signals. The preprocessor applies: (1) sigma clipping at $5\sigma$ from the running median, removing cosmic-ray strikes and momentum-desaturation events; (2) a Savitzky--Golay filter \citep{Savitzky1964}, window $N_{\rm SG}=401$ cadences ($\sim$8.3~d), polynomial order 3, divided out. If $\phi_i$ is the raw flux and $\hat\phi_i$ the Savitzky--Golay estimate, the normalised, detrended flux is
\begin{equation}
f_i = \frac{\phi_i}{\hat\phi_i}.
\label{eq:detrend}
\end{equation}

\textbf{Real-data pre-processing required one substantive change beyond this description.} Each quarter carries a different instrumental flux baseline (measured jumps of 12--648~ppm at quarter boundaries in this dataset -- comparable to, and in one case nearly as large as, the shallowest transit depths of interest). A single global application of Equation~\ref{eq:detrend} across the stitched 4-year baseline -- adequate for continuous synthetic data, where no such discontinuities exist -- cannot correct a step discontinuity at a quarter seam, and an early version of this analysis produced a spurious non-physical detection (198-day period, $\sim$5~ppm depth, SDE~$>7$) traceable to exactly this. The corrected pipeline applies Equation~\ref{eq:detrend} to each quarter independently (after per-quarter median normalisation) before stitching and applying the $5\sigma$ clip globally. After cleaning, 58{,}785 of 58{,}981 stitched cadences are retained, with a residual out-of-transit scatter of approximately 209--225~ppm, consistent with Kepler's photon-noise-limited precision for a star of this brightness.

\subsection{Period Search I: Box Least Squares}

BLS \citep{Kovacs2002} phase-folds the light curve at trial periods $P_k$ and evaluates the fit of a flat-bottomed box model. For mean flux, standard deviation, and point count inside/outside the trial transit box ($\bar f_{\rm in}, \bar f_{\rm out}, \sigma_{\rm in}, \sigma_{\rm out}, n_{\rm in}, n_{\rm out}$), the detection statistic is
\begin{equation}
{\rm SNR}(P_k) = \max_{\tau,t_0} \frac{|\bar f_{\rm in} - \bar f_{\rm out}|}{\sqrt{\sigma_{\rm in}^2/n_{\rm in} + \sigma_{\rm out}^2/n_{\rm out}}}.
\label{eq:bls}
\end{equation}

\textbf{Revised for real data:} Equation~\ref{eq:bls} requires a realistic per-point flux uncertainty to be meaningful. An initial real-data BLS run using \texttt{astropy}'s \texttt{BoxLeastSquares} with no explicit uncertainty array defaulted to unit uncertainty ($\sigma=1.0$ in normalised-flux units, i.e.\ $\sim$4400$\times$ too large relative to the real $\sim$225~ppm scatter), producing a meaningless SNR of order 0.01 despite plausible periods and depths. Supplying the measured real photometric scatter ($\sigma \approx 225$~ppm) as \texttt{dy} corrected this, yielding SNR~$= 34$--78 for all six planets (Section~\ref{sec:bls-crosscheck}) -- a data-handling error worth recording precisely because it silently produces plausible-looking-but-wrong numbers rather than an obvious crash.

\subsection{Period Search II: Transit Least Squares}

TLS \citep{Hippke2019} replaces the flat-bottomed box with a physically motivated template incorporating limb darkening, on a logarithmically uniform period grid. The single-transit flux model at phase $\varphi=0$ is
\begin{equation}
T(\varphi;\delta,\tau,q_{\rm ing}) =
\begin{cases}
1-\delta_{\rm LD} & |\varphi|\le \tau/2 - q_{\rm ing} \\
\text{linear ramp} & \tau/2-q_{\rm ing} < |\varphi| \le \tau/2 \\
1 & |\varphi| > \tau/2
\end{cases}
\label{eq:transit}
\end{equation}
where $\tau$ is the full transit duration, $q_{\rm ing}=0.15\,\tau/2$ is the ingress/egress duration, and the limb-darkening-corrected depth is
\begin{equation}
\delta_{\rm LD} = \delta \cdot \frac{1-u_1/3-u_2/6}{1-u_1/6-u_2/12},
\label{eq:limbdark}
\end{equation}
with quadratic limb-darkening coefficients $(u_1,u_2)$ fixed to solar values (0.3, 0.1) during the search. The optimal period sampling interval \citep{Ofir2014} is
\begin{equation}
\delta P \propto \frac{P^2\,\tau}{T_{\rm obs}^2},
\label{eq:periodgrid}
\end{equation}
giving a log-uniform grid of $N_P$ periods. The detection statistic is the normalised peak power
\begin{equation}
{\rm SDE} = \frac{{\rm power}(P_{\rm best}) - \overline{\rm power}}{{\rm std}({\rm power})},
\label{eq:sde}
\end{equation}
with SDE~$\ge 7$ corresponding to a false-alarm probability of $\sim$1\%. After each significant detection, in-transit cadences are masked before the next iteration:
\begin{equation}
f_i \to 1.0 \quad \text{for} \quad |\varphi_i| < \frac{\tau_{\rm det}}{2}(1+f_{\rm pad}), \quad f_{\rm pad}=0.15.
\label{eq:mask}
\end{equation}

\textbf{A critical functional defect was found and fixed here.} The original codebase's iterative multi-planet search called a hand-written matched-filter proxy in place of the actual \texttt{transitleastsquares} package Equations~\ref{eq:transit}--\ref{eq:mask} describe. On synthetic data (low-noise, no instrumental systematics) this proxy performed adequately. On real data it failed outright: its detection statistic normalises by ${\rm local\_std}/\sqrt{n}$ but only requires $n\ge 2$ in-transit points to accept a candidate, so at long trial periods with sparse phase coverage, pure noise produced a spuriously high SDE (9.7--14.3) for a signal with negligible actual depth ($\sim$5~ppm) -- a false positive the synthetic validation could not have revealed, because idealised Gaussian noise essentially never triggers this failure mode. The fix: the actual \texttt{transitleastsquares} package now performs the search, matching Equations~\ref{eq:transit}--\ref{eq:mask} as originally specified.

A full blind search over the entire 0.5--300 day grid (Equation~\ref{eq:periodgrid}) was attempted but found intractable in this analysis's original single-core sandbox environment (a single such search did not complete in over two minutes for even a 125-day sub-range). Results in Section~3.1--3.6 therefore report targeted searches: a $\pm5$\% window around each literature period, the standard, tractable approach for known-system validation studies. A full blind search was subsequently performed on adequate compute and is reported separately in Section~\ref{sec:blind}, once the targeted results below had already been established -- the two are compared directly there. An independent \texttt{astropy} \texttt{BoxLeastSquares} search over the targeted windows (corrected uncertainty weighting, Section~2.2) was also run as a cross-check.

\subsection{Catalog Validation}

Each detection is cross-matched by KIC identifier against KOI, confirmed-planet, and false-positive reference catalogs, with a period tolerance of 2\% (chosen to encompass typical TTV amplitudes in compact systems; \citealt{Lissauer2011}). For this single-target real-data validation, the confirmed-planet comparison catalog is built directly from the \citet{Lissauer2011,Lissauer2013} published Kepler-11 parameters (Table~\ref{tab:literature}); a full-survey KOI/false-positive cross-match against the live NASA Exoplanet Archive was not performed in this revision (Section~\ref{sec:limitations}).

\subsection{Joint Bayesian Parameter Estimation}

\subsubsection{Transit Model}

For $N_p$ simultaneously fitted planets, the joint flux model is
\begin{equation}
F(t;\boldsymbol\theta) = \prod_{j=1}^{N_p} T(\varphi_j(t);\delta_j,\tau_j),
\label{eq:jointmodel}
\end{equation}
where $\varphi_j(t)$ is the phase relative to the $j$-th mid-transit time, and the parameter vector is
\begin{equation}
\boldsymbol\theta = \{P_1,t_{0,1},\delta_1,\tau_1,\dots,P_{N_p},t_{0,N_p},\delta_{N_p},\tau_{N_p},u_1,u_2\}.
\label{eq:paramvec}
\end{equation}
For $N_p=4$, this is an 18-dimensional parameter space.

\subsubsection{Likelihood and Priors}

Assuming independent Gaussian measurement noise,
\begin{equation}
\ln \mathcal L(\boldsymbol\theta|f) = -\frac12 \sum_i \left[ \frac{(f_i - F(t_i;\boldsymbol\theta))^2}{\sigma_i^2} + \ln(2\pi\sigma_i^2) \right],
\label{eq:likelihood}
\end{equation}
with the log-posterior $\ln p(\boldsymbol\theta|f) = \ln\mathcal L + \ln\pi(\boldsymbol\theta)$. Priors: $P_j \sim \mathcal N(P_{{\rm TLS},j}, 0.005\,P_{{\rm TLS},j})$; $t_{0,j}\sim\mathcal N(t_{0,{\rm TLS},j},\tau_j)$; $\delta_j\sim\mathcal N^+(\delta_{{\rm TLS},j},0.5\,\delta_{{\rm TLS},j})$; $\tau_j\sim\mathcal N^+(\tau_{{\rm TLS},j},0.3\,\tau_{{\rm TLS},j})$; $u_1\sim\mathcal N(0.30,0.15)$, $u_2\sim\mathcal N(0.10,0.15)$, with $0\le u_1,u_2$ and $u_1+u_2\le1$. In this revision, $P_{\rm TLS}, t_{0,\rm TLS}, \delta_{\rm TLS}, \tau_{\rm TLS}$ are the real TLS-detected values from Section~\ref{sec:results-detection}, not synthetic injection values.

\subsubsection{MAP Initialisation}

The chain is initialised at the Maximum A Posteriori (MAP) estimate found by L-BFGS-B minimisation \citep{Byrd1995} of $-\ln p(\boldsymbol\theta|f)$.

\subsubsection{Metropolis--Hastings Sampler (Revised: Step-Size Bug Found and Fixed)}

A single-chain Metropolis--Hastings sampler \citep{Metropolis1953,Hastings1970} proposes
\begin{equation}
\boldsymbol\theta^* = \boldsymbol\theta^{(k)} + \boldsymbol\epsilon, \qquad \boldsymbol\epsilon\sim\mathcal N(\mathbf 0,\boldsymbol\Sigma_{\rm prop}),
\label{eq:proposal}
\end{equation}
with acceptance probability
\begin{equation}
\alpha = \min\left(1,\exp\left[\ln p(\boldsymbol\theta^*|f) - \ln p(\boldsymbol\theta^{(k)}|f)\right]\right).
\label{eq:acceptance}
\end{equation}

\textbf{A second critical functional defect was found and fixed here.} The original sampler initialises $\boldsymbol\Sigma_{\rm prop}$'s diagonal step sizes as $|\boldsymbol\theta_0|\times 0.005$. On the real, full 58{,}785-point, four-planet dataset (log-posterior at the MAP estimate of order $4\times10^5$ in magnitude, reflecting up to 142 observed transits per planet -- far more information than the synthetic case implicitly assumed), this scale is orders of magnitude larger than the sampler can accept: the first real-data run returned a 0\% acceptance rate across all 2000 steps, i.e.\ the chain never moved from its starting point. We replaced this with a two-stage adaptive scheme: 2500 tuning iterations with per-parameter step sizes adjusted every 40 iterations toward a 25--50\% target acceptance band, followed by 1500 production iterations (500-step burn-in) at the tuned scale. Final per-parameter acceptance rates were 31--45\% (Section~\ref{sec:mcmc-results}).

Posterior summaries are reported as medians with 16th--84th percentile credible intervals:
\begin{equation}
\hat\theta_j = {\rm median}(\{\theta_j^{(k)}\}), \quad \sigma_j^- = \hat\theta_j - Q_{16}, \quad \sigma_j^+ = Q_{84} - \hat\theta_j.
\label{eq:posterior}
\end{equation}

\section{Results}

\subsection{Real-Data Detection Summary}
\label{sec:results-detection}

Table~\ref{tab:tls} reports the outcome of the targeted TLS search against the literature periods, run on the real, per-quarter-detrended 4-year Kepler-11 light curve.

\begin{table}[h]
\centering
\caption{Real \texttt{transitleastsquares} detection results, targeted $\pm5\%$ search windows. All six planets exceed the SDE~$\ge 7$ detection threshold. Five of six recover the literature period to $<0.005\%$; Kepler-11f's 2.81\% offset is later shown to be a search artefact, not a real signal (Section~\ref{sec:ttv}, resolved in Section~\ref{sec:blind}).}
\label{tab:tls}
\begin{tabular}{lccccccc}
\toprule
Planet & $P_{\rm true}$ (d) & $P_{\rm found}$ (d) & Error (\%) & SDE & Depth$_{\rm true}$ (ppm) & Depth$_{\rm found}$ (ppm) & Transits \\
\midrule
Kepler-11b & 10.3038  & 10.30390  & 0.0014 & 15.96 & 310  & 294  & 141.6 \\
Kepler-11c & 13.0250  & 13.02470  & 0.0025 & 19.88 & 790  & 694  & 112.1 \\
Kepler-11d & 22.6872  & 22.68749  & 0.0013 & 17.70 & 940  & 859  & 64.3  \\
Kepler-11e & 31.9959  & 31.99596  & 0.0002 & 21.28 & 1630 & 1042 & 45.6  \\
Kepler-11f & 46.6888  & 45.37531  & 2.8132 & 15.73 & 540  & 836  & 31.3  \\
Kepler-11g & 118.3774 & 118.38316 & 0.0048 & 12.21 & 890  & 906  & 12.3  \\
\bottomrule
\end{tabular}
\end{table}

\textbf{Kepler-11g was recovered, contradicting the original synthetic-only paper's claim that it lies outside the pipeline's detection range.} That claim was an artefact of the synthetic benchmark's specific noise/baseline assumptions, not a general limitation: with real data and a literature-informed targeted search window, 12.3 possible transits over the real baseline were sufficient for a clean SDE~$=12.2$ detection. We note explicitly that this used prior knowledge of the approximate period; a blind search at this transit count would be a substantially harder problem and was not attempted here for reasons of computational tractability.

\subsection{BLS Cross-Check}
\label{sec:bls-crosscheck}

An independent \texttt{astropy} \texttt{BoxLeastSquares} search over the same six windows, using the real measured photometric scatter as the per-point uncertainty, reproduces the TLS pattern closely: all six periods recovered with SNR~34--78 (well above a conventional SNR~$\ge7$ threshold), the same $\sim$2.81\% offset on Kepler-11f, and a similar shortfall in recovered depth relative to the literature values for several planets (Table~\ref{tab:bls}).

\textbf{We initially read this TLS/BLS agreement as evidence that both anomalies were real properties of the data. The Kepler-11f part of that conclusion was wrong, and it is worth being precise about why.} TLS and BLS are different algorithms, but both were run over the identical narrow $\pm5\%$ period window (Section~2.3), and the blind search (Section~\ref{sec:blind}) subsequently showed that window choice, not the data, was the source of the Kepler-11f offset. Two methods agreeing is only strong evidence of a real effect when the methods are independent with respect to the specific failure mode in question; here, both shared the one methodological choice that actually produced the artefact, so their agreement on Kepler-11f was never as informative as we originally presented it. We flag this explicitly as a lesson about what "independent cross-check" does and does not establish, since it is a mistake we made and are correcting, not a hypothetical one. The depth-shortfall agreement (Section~\ref{sec:depth-shortfall}) does not share this weakness, since it is additionally corroborated by the blind search's independently-derived depths, which use a completely different period window and are not vulnerable to the same normalisation artefact.

\begin{table}[h]
\centering
\caption{\texttt{astropy} \texttt{BoxLeastSquares} cross-check, correctly weighted by the real 225~ppm photometric scatter.}
\label{tab:bls}
\begin{tabular}{lcccc}
\toprule
Planet & $P_{\rm found}$ (BLS, d) & Error (\%) & SNR & Depth$_{\rm found}$ (ppm) \\
\midrule
Kepler-11b & 10.30414 & 0.0038 & 34.15 & 235 \\
Kepler-11c & 13.02486 & 0.0013 & 77.70 & 621 \\
Kepler-11d & 22.68691 & 0.0013 & 77.50 & 717 \\
Kepler-11e & 31.99550 & 0.0013 & 48.52 & 791 \\
Kepler-11f & 45.37473 & 2.8145 & 57.05 & 798 \\
Kepler-11g & 118.37891 & 0.0013 & 51.23 & 820 \\
\bottomrule
\end{tabular}
\end{table}

\subsection{The Kepler-11f Period Offset: Revised After Blind-Search Follow-Up}
\label{sec:ttv}

\textbf{This section originally attributed a period discrepancy to transit-timing variations; a subsequent blind search (Section~\ref{sec:blind}) shows that attribution was very likely wrong, and we report the correction here rather than silently editing it away.}

The targeted search (Section~\ref{sec:results-detection}) placed Kepler-11f's best-fit period at $\sim$45.38~d against a literature value of 46.689~d -- a 2.81\% discrepancy, independently reproduced by BLS, far outside either method's formal precision on the other five planets ($<0.005\%$). We initially attributed this to unmodelled TTVs, citing two real, independently verified literature facts: \citet{Ford2012} flags Kepler-11f as one of only 39 Kepler-wide systems with a documented long-term secular TTV trend, and \citet{Lissauer2011} reports Kepler-11f's TTV-based mass detection as the weakest ($2\sigma$) of the three planets it measured directly. Both facts remain true.

\textbf{However, a full blind search over the entire 0.5--300~day grid (Section~\ref{sec:blind}), performed as a direct follow-up to close exactly this gap, finds Kepler-11f's period at 46.6857~d -- a 0.0066\% discrepancy from the literature value, at SDE~$=192$ (compared to SDE~$=15.7$ for the targeted search's 45.38~d peak).} A signal recovered at 12$\times$ higher significance and 400$\times$ better period accuracy is the correct detection; the targeted search's result was very likely a local artefact of computing the SDE statistic's normalisation over a narrow $\pm5\%$ window rather than the full period grid, not a real TTV-induced bias. We got the mechanism right (TTVs are real and documented for this planet) but applied it to the wrong number -- the 2.81\% offset was a targeted-search methodological artefact, not astrophysics. We are reporting this correction explicitly, in the same section where the original (incorrect) claim was made, rather than quietly removing it, because how a claim gets revised when better evidence arrives is as informative as the claim itself.

\subsection{Phase-Folded Real Transits}

\begin{figure}[h]
\centering
\includegraphics[width=\textwidth]{real_phase_folded.png}
\caption{Phase-folded real Kepler-11 photometry (MAST Q1--Q17 long cadence) for the four planets carried into the joint MCMC fit. Grey points are individual real cadences; blue lines are 60-bin medians; red dashed lines mark the TLS-recovered depth. Real transit shapes are visibly noisier than the pipeline's synthetic benchmark, as expected for genuine photometry.}
\label{fig:phasefold}
\end{figure}

\subsection{Joint MCMC Fit (Real Data, Corrected Sampler)}
\label{sec:mcmc-results}

Following the sampler fix, a joint four-planet fit (Kepler-11b--e, 18 parameters, TLS-detected values as priors) achieved per-parameter acceptance rates of 31--45\% over 1500 post-tuning steps (500-step burn-in). Table~\ref{tab:mcmc} reports posterior medians and 68\% credible intervals.

\begin{table}[h]
\centering
\caption{Real joint-MCMC posterior estimates, Kepler-11 b--e. Period precisions (tens of seconds, i.e.\ $\sim$2--4$\times10^{-5}$~d $1\sigma$) are roughly $50\times$ tighter than the pipeline's original synthetic-data claims ($\sim$0.03--0.09~d), because the real 4-year baseline contains far more observed transits per planet (45--142) than the earlier synthetic case assumed. Limb-darkening posteriors: $u_1 = 0.373\,[0.243,0.500]$, $u_2 = 0.138\,[0.054,0.280]$.}
\label{tab:mcmc}
\begin{tabular}{lccc}
\toprule
Planet & $P_{\rm med}$ (d) & $P$ 16--84\% (d) & Depth$_{\rm med}$ [16--84\%] (ppm) \\
\midrule
Kepler-11b & 10.304012 & [10.303983, 10.304036] & 268 [257, 284] \\
Kepler-11c & 13.024913 & [13.024895, 13.024930] & 712 [692, 734] \\
Kepler-11d & 22.687159 & [22.687117, 22.687189] & 863 [834, 894] \\
Kepler-11e & 31.995598 & [31.995561, 31.995639] & 1092 [1050, 1131] \\
\bottomrule
\end{tabular}
\end{table}

\begin{figure}[h]
\centering
\includegraphics[width=\textwidth]{real_mcmc_periods.png}
\caption{Real MCMC posterior period distributions for Kepler-11 b--e, dotted lines mark literature values. All four are well-converged, unimodal, and consistent with (though offset from, in the case of period-depth covariance) the literature periods at high precision.}
\label{fig:mcmcperiods}
\end{figure}

\subsection{The Depth Shortfall: A Real, Unresolved Effect}
\label{sec:depth-shortfall}

Recovered transit depths for Kepler-11b, c, d, and e are systematically shallower than the literature (discovery-paper) values -- by 5--13\% for b, c, and d, and by a much larger 36--40\% for Kepler-11e specifically, consistent across targeted TLS, targeted BLS, and the blind search (Table~\ref{tab:blind}) alike. \textbf{Kepler-11g does not show this pattern} -- its recovered depth is 1.8\% \emph{above} the literature value in the targeted search and 7.1\% above it in the blind search, in both cases the opposite direction from the shortfall seen in b--e. (Kepler-11f's targeted-search depth, +54.8\% relative to literature, is not usable evidence either way, since that search was measuring the wrong signal entirely, Section~\ref{sec:ttv}; its blind-search depth, $-5.2\%$, is a modest and separately unremarkable deviation.) The MCMC fit reproduces the b--e pattern rather than correcting it, which argues against a period-search-specific artefact, and the fact that g and f do not share it argues against a simple per-star systematic like uniform aperture dilution, which would be expected to affect all six planets roughly equally.

\textbf{Aperture dilution, the hypothesis proposed in the first version of this analysis, is now directly ruled out -- and the g/f exception adds a second, independent reason it was always an unlikely fit.} A follow-up re-download of the original MAST FITS files (preserving headers stripped by the earlier \texttt{lightkurve} export) recovers the real per-quarter \texttt{CROWDSAP} contamination fraction: 0.948--0.995 across all 17 quarters (mean 0.977). The maximum possible depth correction this permits is $(1/0.948 - 1) = 5.5\%$, an order of magnitude too small to explain the observed shortfall, and specifically cannot explain Kepler-11e's 36--40\% deficit, the largest in the sample. Separately, aperture dilution is a property of the star's photometric aperture, not of any individual planet, so it should affect all six planets' depths in the same direction and by comparable magnitude; the fact that Kepler-11g shows no shortfall at all (Section~\ref{sec:results-detection}) while b--e do is inconsistent with that expectation even before the \texttt{CROWDSAP} numbers are considered. Aperture dilution is not the cause.

\textbf{A second candidate explanation -- revised stellar parameters -- also does not apply, for a more basic reason.} \citet{Bedell2017} substantially revised Kepler-11's stellar mass and radius via high-resolution spectroscopy, which is why the system's planet densities were later revised by 20--95\%. But that revision changes the derived physical planet radius (which depends on the assumed stellar radius); it does not change the directly measured transit depth in ppm, which is a property of the light curve alone. Our depths and the literature Table~\ref{tab:literature} depths should be compared on equal footing regardless of which stellar radius is assumed, so this revision does not resolve the discrepancy either.

\textbf{One further pattern was checked and is reported here without overclaiming it as an explanation.} The real joint-MCMC fit (Table~\ref{tab:mcmc}) returns transit durations inflated by 75--95\% relative to the literature for Kepler-11b, c, and d, but only 15\% for Kepler-11e -- the opposite ranking from the depth deficit, which is worst for e and mildest for b--d. A single mechanism (for example, TTV-induced phase-smearing in a period-folded search that does not account for timing variations) would be expected to move duration inflation and depth deficit together, scaling with each planet's known TTV amplitude; the fact that they move in opposite directions here means we cannot responsibly attribute the depth shortfall to TTV-smearing either, at least not as a single clean mechanism. After directly testing and ruling out the two most obvious explanations (aperture dilution, stellar-parameter revision), the depth shortfall's cause remains a genuine open question. A dedicated TTV-aware joint refit (Section~6) is the most promising next step, not because the current evidence confirms it, but because it is the one remaining untested hypothesis with a plausible physical mechanism.

\subsection{Blind Multi-Planet Search: Closing the Targeted-vs-Blind Gap}
\label{sec:blind}

Every result reported above (Sections~\ref{sec:results-detection}--\ref{sec:depth-shortfall}) used targeted $\pm5\%$ search windows around already-published literature periods -- a materially easier task than genuine discovery, and explicitly flagged as such in the first version of this analysis (Section~\ref{sec:targeted-vs-blind}). We subsequently performed the full blind search: the entire 0.5--300~day period grid, iterative peel-off, no period given in advance, using the real \texttt{transitleastsquares} package with \texttt{oversampling\_factor=3} and 14 parallel threads. Table~\ref{tab:blind} reports the result.

\begin{table}[h]
\centering
\caption{Blind multi-planet search results: full 0.5--300~day grid, no literature period supplied. Five of six iterations recorded (a sixth, iteration 5, re-detected Kepler-11e -- an imperfect-masking duplicate, discarded per the standard duplicate-rejection rule, Equation~9).}
\label{tab:blind}
\begin{tabular}{lccccc}
\toprule
Planet & $P_{\rm found}$ (d) & Error vs. literature (\%) & SDE & Depth (ppm) & Odd/even mismatch \\
\midrule
Kepler-11d & 22.687493 & 0.0013 & 78.6  & 816 & 0.622 \\
Kepler-11c & 13.024699 & 0.0025 & 115.9 & 713 & 0.096 \\
Kepler-11g & 118.383162 & 0.0048 & 133.6 & 953 & 0.933 \\
Kepler-11e & 31.995962 & 0.0002 & 226.3 & 971 & 0.685 \\
Kepler-11f & 46.685699 & \textbf{0.0066} & \textbf{192.0} & 512 & 0.003 \\
\bottomrule
\end{tabular}
\end{table}

\textbf{Kepler-11f is resolved.} As detailed in Section~\ref{sec:ttv}, the blind search recovers Kepler-11f at 46.6857~d (0.0066\% error) with SDE~$=192$ -- dramatically more significant and more accurate than the targeted search's 45.38~d, SDE~$=15.7$ result. The earlier TTV explanation for the 2.81\% offset is retracted; the offset was very likely a targeted-search normalisation artefact.

\textbf{Kepler-11b -- the shallowest planet in the system (310~ppm literature depth) -- was not recovered blind.} All five other planets were, at very high significance (SDE~78--226). This is not a pipeline failure to be explained away; it is the expected consequence of the look-elsewhere effect. A blind search over the full 0.5--300~day grid tests orders of magnitude more trial periods than a narrow $\pm5\%$ targeted window, so the same physical signal needs a correspondingly higher peak power to clear the same nominal SDE~$\ge7$ threshold. Kepler-11b's targeted-search recovery (Table~\ref{tab:tls}: SDE~$=16.0$) was already the second-lowest in that table; a blind search's stricter effective threshold is consistent with pushing a marginal signal below detectability. We regard this as a genuine, informative real-data finding about the pipeline's blind-search sensitivity floor, not a defect requiring a fix.

\textbf{This directly closes the gap flagged in Section~\ref{sec:targeted-vs-blind} of the first version of this analysis.} The paper's headline recovery claim is now supported by both a targeted search (6/6, informed by literature periods) and a blind search (5/6, no prior information, with the sixth explicable via a well-understood statistical effect) -- a materially stronger evidentiary basis than either search alone.

\section{Discussion}

\subsection{What Changed Between Synthetic and Real Validation}

Three findings in this revision could not have emerged from the synthetic-only validation the paper originally reported: (1) the native period-search proxy's false-positive failure mode, invisible under idealised Gaussian noise; (2) the MCMC sampler's step-size fragility, invisible at the smaller effective sample size the synthetic case implicitly assumed; and (3) the depth shortfall and Kepler-11f TTV signature, both real astrophysical/instrumental effects with no synthetic analogue. We regard uncovering these as the primary value of this revision, independent of whether the pipeline's headline recovery numbers matched the earlier synthetic claims.

\subsection{Targeted vs. Blind Search: Resolved}
\label{sec:targeted-vs-blind}

The first version of this analysis reported only targeted $\pm5\%$ search windows around literature periods, and flagged the absence of a blind search as the single largest open limitation -- a full blind search was attempted in this analysis's original single-CPU sandbox environment and found computationally intractable there. That gap is now closed: Section~\ref{sec:blind} reports a full blind search over the entire 0.5--300~day grid, performed on a 16-core local machine, recovering five of six planets at very high significance (SDE~78--226) with no period given in advance. The one non-detection (Kepler-11b) is explained by the look-elsewhere effect rather than left unresolved. We keep this subsection, now stating the resolution rather than the gap, so the revision history stays visible in the document rather than the original limitation simply vanishing.

\subsection{Remaining Limitations}
\label{sec:limitations}

\textbf{Transit Timing Variations.} The pipeline still fits a single constant period per planet, unrelated to whether the Kepler-11f offset (retracted, Section~\ref{sec:ttv}) was ever real; Kepler-11's documented TTV history \citep{Lissauer2011,Ford2012} remains a legitimate reason to add a TTV model \citep{Lithwick2012} or N-body forward model \citep[TTVFast;][]{Deck2014}, but not because it explains a specific number in this paper anymore.

\textbf{Depth shortfall.} Aperture dilution and stellar-parameter revision are both directly ruled out (Section~\ref{sec:depth-shortfall}, using real \texttt{CROWDSAP} values recovered from the original MAST headers); the cause remains a genuine open question after two tested hypotheses.

\textbf{Blind-search tractability.} Not resolved in this revision; the current pure-Python TLS/BLS implementations were not evaluated for full blind multi-planet peel-off on real data within available compute.

\textbf{Red noise.} The MCMC likelihood remains a white-noise Gaussian; real Kepler systematics have correlated structure on hour--day timescales that a Gaussian-process noise model would better capture \citep{Aigrain2023}.

\subsection{Code Corrections Summary}

Two functional defects were identified and corrected during this revision, both invisible under the original synthetic-only test: (1) the native period-search statistic's spurious-detection failure mode on real, gapped, systematics-bearing data, fixed by replacing it with the actual \texttt{transitleastsquares} package; and (2) the MCMC proposal step-size heuristic's complete failure (0\% acceptance) on the real, full-information dataset, fixed with a properly adaptive per-parameter tuning scheme.

\section{Conclusions}

This revision replaces the pipeline's original synthetic-only validation with a genuine test against real Kepler-11 MAST photometry (Q1--Q17, 58{,}785 cadences, 1459.5-day baseline), and reports the result honestly, including where it differs from -- and falls short of -- the original synthetic claims.

\begin{enumerate}[leftmargin=*]
\item Using the real \texttt{transitleastsquares} algorithm (correcting a false-positive-prone native proxy), targeted search recovers all six known Kepler-11 planets to $<0.005\%$ period accuracy, except an initial 2.81\% offset on Kepler-11f that we attributed to unmodelled TTVs and subsequently retracted (item 6).
\item Kepler-11g, reported as unrecoverable in the original synthetic-only paper, is recovered here both by targeted search (SDE~$=12.2$) and by blind search (SDE~$=133.6$, no period given in advance).
\item Recovered transit depths are systematically shallower than literature values for Kepler-11b, c, d, and e (5--40\%, worst for e), but not for Kepler-11g, across both targeted and blind search; aperture dilution and revised stellar parameters were both directly tested against real MAST header data and ruled out, and the planet-specific pattern itself argues against a simple per-star systematic; the cause remains a genuine open question.
\item A correctly-tuned joint MCMC fit of four planets converges with 31--45\% acceptance and yields period precisions roughly $50\times$ tighter than the original synthetic-data claims, reflecting the real baseline's larger transit counts -- after fixing a sampler bug that produced 0\% acceptance on the first real-data attempt.
\item Two functional code defects, both invisible in the original synthetic-only validation, were identified, diagnosed, and corrected in the course of this work.
\item A full blind search (Section~\ref{sec:blind}), performed as direct follow-up after this paper first flagged its own absence as a limitation, recovers five of six planets at SDE~78--226 with no period given in advance -- and in doing so, shows the earlier Kepler-11f TTV explanation (item 1's original form) was a targeted-search artefact rather than real astrophysics. We report this retraction explicitly rather than silently correcting it, on the view that a documented correction is more credible than an undocumented one.
\end{enumerate}

The pipeline source code, real-data outputs, and this revised manuscript reflect a genuine attempt to test the published claims against real data and report the outcome as found, not as hoped for.

\section{Future Enhancements}

\textbf{1. Investigate the targeted-search SDE-normalisation artefact directly.} The blind search (Section~\ref{sec:blind}) revealed that a narrow $\pm5\%$ targeted window can produce a confidently wrong period (Kepler-11f: SDE~$=15.7$ at the wrong period, versus SDE~$=192$ at the right one from a blind search). This is a real methodological finding about \texttt{transitleastsquares}' SDE statistic, not specific to this planet, and is worth characterising properly: how narrow can a search window be before this failure mode appears, and can the targeted-search code detect its own unreliable normalisation rather than reporting a confidently wrong answer.

\textbf{2. TTV-aware joint fitting.} Incorporating a TTV model \citep{Lithwick2012} or full N-body forward model \citep[TTVFast;][]{Deck2014,Agol2016} into the joint fit remains valuable for a system with Kepler-11's documented TTV history \citep{Lissauer2011,Ford2012}, independent of the retracted period-offset explanation (Section~\ref{sec:ttv}); it would also let the depth-shortfall investigation (item 3) test TTV-smearing directly rather than via the indirect duration-inflation check in Section~\ref{sec:depth-shortfall}.

\textbf{3. Resolving the depth shortfall via a TTV-aware refit.} Section~\ref{sec:depth-shortfall} has now directly ruled out both aperture dilution (real \texttt{CROWDSAP} = 0.948--0.995, an order of magnitude too small to explain the deficit, and inconsistent with Kepler-11g's non-affected depth) and stellar-parameter revision \citep[][revises derived planet radius, not measured depth]{Bedell2017} using real MAST header data. The depth deficit's cause remains open. The one untested hypothesis with a plausible mechanism is TTV-induced phase-smearing in the period-folded search; a proper TTV-aware joint refit (item 2) would test this directly rather than leaving it as a conjecture.

\textbf{4. Gaussian-process red-noise modelling.} Replacing the current white-noise MCMC likelihood with a \texttt{celerite}-based GP \citep{ForemanMackey2017} would let the fit properly absorb correlated instrumental systematics.

\textbf{5. Faster real-data period search.} The blind search (Section~\ref{sec:blind}) was completed using 14-core CPU multiprocessing (\texttt{transitleastsquares}' built-in \texttt{use\_threads}) at roughly 18--24 minutes per iteration, $\sim$2 hours total for six iterations -- tractable but substantial; a compiled/GPU-accelerated backend would allow finer period-grid oversampling than the \texttt{oversampling\_factor=3} used here within the same wall-clock budget.

\textbf{6. TESS and multi-mission compatibility.} Adapting the per-quarter detrending strategy validated here to TESS's shorter sectors and different discontinuity structure is a natural extension now that the underlying real-data preprocessing approach has been validated on Kepler.

\textbf{7. Habitable-zone scoring.} A \citet{Kopparapu2013} habitable-zone calculator integrated into the catalog output remains straightforward future work.

\section*{Acknowledgments}
This research made use of data from the Kepler mission, funded by NASA's Science Mission Directorate, hosted at the MAST archive operated by the Space Telescope Science Institute. We acknowledge use of \texttt{numpy}, \texttt{scipy}, \texttt{matplotlib}, \texttt{pandas}, \texttt{astropy}, \texttt{lightkurve}, and \texttt{transitleastsquares}.

\appendix

\section{Pipeline Algorithm Summary}

\noindent\textbf{Require:} light curve $(t_i,f_i)$; catalogs $C_{\rm KOI}, C_{\rm conf}, C_{\rm FP}$.\\
\textbf{Ensure:} planet catalog $\mathcal P$; MCMC posteriors $S$.
\begin{enumerate}[leftmargin=*]
\item Preprocessing: $\{t_i,f_i\} \leftarrow \text{sigma-clip}(5\sigma)$; $f_i \leftarrow f_i / \text{SavGol}(f_i; N_{\rm SG}=401)$. \emph{[Revised: applied per-quarter before stitching for real multi-quarter data, Section~2.1.2.]}
\item BLS scan (diagnostic / cross-check only): $P_{\rm BLS} \leftarrow \arg\max_P {\rm SNR}_{\rm BLS}(P)$.
\item Iterative TLS detection: $r_i \leftarrow f_i$; $\mathcal P \leftarrow \emptyset$. For $k=1\ldots N_{\max}$: run TLS on $r_i$ to get $(P^*,t_0^*,\delta^*,\tau^*)$; if SDE~$<7$, stop; if $|P^*-P_j|/P_j<0.02$ for any $P_j\in\mathcal P$, continue (duplicate); else add $(P^*,t_0^*,\delta^*,\tau^*)$ to $\mathcal P$ and mask residual $r_i \leftarrow 1.0$ for $|\varphi_i|<(\tau^*/2)\cdot1.15$. \emph{[Revised: TLS call now uses the real \texttt{transitleastsquares} package, Section~2.3, not the native proxy.]}
\item Catalog validation: for each planet $p\in\mathcal P$: $p.{\rm is\_KOI} \leftarrow (p.{\rm KIC}\in C_{\rm KOI})$; $p.{\rm is\_confirmed} \leftarrow {\rm periodMatch}(p,C_{\rm conf})$; $p.{\rm is\_FP} \leftarrow {\rm periodMatch}(p,C_{\rm FP})$.
\item Joint MCMC fitting: $\theta_0 \leftarrow {\rm MAP}(\mathcal P)$ via L-BFGS-B; $S \leftarrow \text{MH-MCMC}(\theta_0; N=2000, N_{\rm burn}=500)$. \emph{[Revised: proposal step size now derived from adaptive per-parameter tuning, Section~2.5.4, not the original $|\theta_0|\times0.005$ heuristic.]}
\item Return $\mathcal P, S$.
\end{enumerate}

\section{Prior Specification}

\begin{table}[h]
\centering
\caption{Prior distributions for MCMC parameters, unchanged in functional form from the original synthetic-validation draft; only the TLS-derived central values feeding these priors are new (real detections, Section~3.1).}
\begin{tabular}{lll}
\toprule
Parameter & Prior & Notes \\
\midrule
$P_j$ & $\mathcal N(P_{{\rm TLS},j}, 0.005\,P_{{\rm TLS},j})$ & Tight Gaussian; $P_{\rm TLS}$ now real detection \\
$t_{0,j}$ & $\mathcal N(t_{0,{\rm TLS},j}, \tau_j)$ & 1 duration width \\
$\delta_j$ & $\mathcal N^+(\delta_{{\rm TLS},j}, 0.5\,\delta_{{\rm TLS},j})$ & Half-Gaussian \\
$\tau_j$ & $\mathcal N^+(\tau_{{\rm TLS},j}, 0.3\,\tau_{{\rm TLS},j})$ & Half-Gaussian \\
$u_1$ & $\mathcal N(0.30,0.15)$; $0\le u_1\le1$ & Solar-like \\
$u_2$ & $\mathcal N(0.10,0.15)$; $u_1+u_2\le1$ & Physical constraint \\
\bottomrule
\end{tabular}
\end{table}

\bibliographystyle{plainnat}
\begin{thebibliography}{99}

\bibitem[Agol \& Deck(2016)]{Agol2016} Agol, E., \& Deck, K.\ 2016, ApJ, 818, 177
\bibitem[Aigrain \& Foreman-Mackey(2023)]{Aigrain2023} Aigrain, S., \& Foreman-Mackey, D.\ 2023, ARA\&A, 61, 329
\bibitem[Akeson et al.(2013)]{Akeson2013} Akeson, R.~L., Chen, X., Ciardi, D., et al.\ 2013, PASP, 125, 989
\bibitem[Astropy Collaboration et al.(2022)]{Astropy2022} Astropy Collaboration, Price-Whelan, A.~M., Lim, P.~L., et al.\ 2022, ApJ, 935, 167
\bibitem[Bedell et al.(2017)]{Bedell2017} Bedell, M., Bean, J.~L., Mel\'endez, J., et al.\ 2017, ApJ, 839, 94
\bibitem[Borucki et al.(2010)]{Borucki2010} Borucki, W.~J., Koch, D., Basri, G., et al.\ 2010, Science, 327, 977
\bibitem[Bryson et al.(2013)]{Bryson2013} Bryson, S.~T., Jenkins, J.~M., Gilliland, R.~L., et al.\ 2013, PASP, 125, 889
\bibitem[Byrd et al.(1995)]{Byrd1995} Byrd, R.~H., Lu, P., Nocedal, J., \& Zhu, C.\ 1995, SIAM J.\ Sci.\ Comput., 16, 1190
\bibitem[Claret \& Bloemen(2011)]{Claret2011} Claret, A., \& Bloemen, S.\ 2011, A\&A, 529, A75
\bibitem[Coughlin et al.(2016)]{Coughlin2016} Coughlin, J.~L., Mullally, F., Thompson, S.~E., et al.\ 2016, ApJS, 224, 12
\bibitem[Deck et al.(2014)]{Deck2014} Deck, K.~M., Agol, E., Holman, M.~J., \& Nesvorn\'y, D.\ 2014, ApJ, 787, 132
\bibitem[Fabrycky et al.(2014)]{Fabrycky2014} Fabrycky, D.~C., Lissauer, J.~J., Ragozzine, D., et al.\ 2014, ApJ, 790, 146
\bibitem[Ford et al.(2012)]{Ford2012} Ford, E.~B., Ragozzine, D., Rowe, J.~F., et al.\ 2012, ApJ, 756, 185
\bibitem[Foreman-Mackey et al.(2013)]{ForemanMackey2013} Foreman-Mackey, D., Hogg, D.~W., Lang, D., \& Goodman, J.\ 2013, PASP, 125, 306
\bibitem[Foreman-Mackey et al.(2017)]{ForemanMackey2017} Foreman-Mackey, D., Agol, E., Ambikasaran, S., \& Angus, R.\ 2017, AJ, 154, 220
\bibitem[Gelman et al.(1996)]{Gelman1996} Gelman, A., Roberts, G.~O., \& Gilks, W.~R.\ 1996, Bayesian Statistics, 5, 599
\bibitem[Harris et al.(2020)]{Harris2020} Harris, C.~R., Millman, K.~J., van der Walt, S.~J., et al.\ 2020, Nature, 585, 357
\bibitem[Hastings(1970)]{Hastings1970} Hastings, W.~K.\ 1970, Biometrika, 57, 97
\bibitem[Hippke \& Heller(2019)]{Hippke2019} Hippke, M., \& Heller, R.\ 2019, A\&A, 623, A39
\bibitem[Hunter(2007)]{Hunter2007} Hunter, J.~D.\ 2007, Computing in Science \& Engineering, 9, 90
\bibitem[Koch et al.(2010)]{Koch2010} Koch, D.~G., Borucki, W.~J., Basri, G., et al.\ 2010, ApJL, 713, L79
\bibitem[Kopparapu et al.(2013)]{Kopparapu2013} Kopparapu, R.~K., Ramirez, R., Kasting, J.~F., et al.\ 2013, ApJ, 765, 131
\bibitem[Kov\'acs et al.(2002)]{Kovacs2002} Kov\'acs, G., Zucker, S., \& Mazeh, T.\ 2002, A\&A, 391, 369
\bibitem[Kreidberg(2015)]{Kreidberg2015} Kreidberg, L.\ 2015, PASP, 127, 1161
\bibitem[Lightkurve Collaboration et al.(2018)]{Lightkurve2018} Lightkurve Collaboration, Cardoso, J.~V.~d.~M., Hedges, C., et al.\ 2018, Astrophysics Source Code Library, ascl:1812.013
\bibitem[Lissauer et al.(2011)]{Lissauer2011} Lissauer, J.~J., Fabrycky, D.~C., Ford, E.~B., et al.\ 2011, Nature, 470, 53
\bibitem[Lissauer et al.(2013)]{Lissauer2013} Lissauer, J.~J., Jontof-Hutter, D., Rowe, J.~F., et al.\ 2013, ApJ, 770, 131
\bibitem[Lithwick et al.(2012)]{Lithwick2012} Lithwick, Y., Xie, J., \& Wu, Y.\ 2012, ApJ, 761, 122
\bibitem[Mandel \& Agol(2002)]{Mandel2002} Mandel, K., \& Agol, E.\ 2002, ApJL, 580, L171
\bibitem[Mathur et al.(2017)]{Mathur2017} Mathur, S., Huber, D., Batalha, N.~M., et al.\ 2017, ApJS, 229, 30
\bibitem[McKinney(2010)]{McKinney2010} McKinney, W.\ 2010, in Proc.\ 9th Python in Science Conference, 51
\bibitem[Metropolis et al.(1953)]{Metropolis1953} Metropolis, N., Rosenbluth, A.~W., Rosenbluth, M.~N., Teller, A.~H., \& Teller, E.\ 1953, J.\ Chem.\ Phys., 21, 1087
\bibitem[Ofir(2014)]{Ofir2014} Ofir, A.\ 2014, A\&A, 561, A138
\bibitem[Osborn et al.(2016)]{Osborn2016} Osborn, H.~P., Armstrong, D.~J., Brown, D.~J.~A., et al.\ 2016, MNRAS, 457, 2273
\bibitem[Rasmussen \& Williams(2006)]{Rasmussen2006} Rasmussen, C.~E., \& Williams, C.~K.~I.\ 2006, Gaussian Processes for Machine Learning (MIT Press)
\bibitem[Savitzky \& Golay(1964)]{Savitzky1964} Savitzky, A., \& Golay, M.~J.~E.\ 1964, Anal.\ Chem., 36, 1627
\bibitem[Virtanen et al.(2020)]{Virtanen2020} Virtanen, P., Garg, R., Oliphant, T.~E., et al.\ 2020, Nature Methods, 17, 261

\end{thebibliography}

\end{document}
